\documentclass[11pt]{article}

% -----------------------------------------------------------------------
%  Packages
% -----------------------------------------------------------------------
\usepackage[margin=1in]{geometry}
\usepackage{booktabs}
\usepackage{amsmath,amssymb}
\usepackage{listings}
\usepackage{xcolor}
\usepackage[T1]{fontenc}
\usepackage[utf8]{inputenc}
\usepackage{hyperref}
\usepackage{graphicx}
\usepackage{multirow}
\usepackage{array}
\usepackage{microtype}
\usepackage{natbib}

\lstset{
  basicstyle=\ttfamily\small,
  breaklines=true,
  frame=single,
  backgroundcolor=\color{gray!5}
}

% -----------------------------------------------------------------------
%  Convenience macros
% -----------------------------------------------------------------------
\newcommand{\cb}{\textsc{ContinuumBench}}
\newcommand{\todo}[1]{\textcolor{red}{\textbf{[TODO: #1]}}}
\newcommand{\RIS}{\mathrm{RIS}}
\newcommand{\viewname}[1]{\texttt{#1}}

% -----------------------------------------------------------------------
%  Title
% -----------------------------------------------------------------------
\title{%
  \cb{}: A Tri-View Benchmark for Representation Invariance\\
  in Relational Machine Learning%
}
\author{%
  Raj Bhalwankar \quad Pasquale Minervini \quad Erman Acar\\[4pt]
  \small \textit{[Institution]}%
}
\date{\today}

% -----------------------------------------------------------------------
\begin{document}
\maketitle

% -----------------------------------------------------------------------
\begin{abstract}
% -----------------------------------------------------------------------
Relational machine learning practitioners must choose how to represent
a database before training: as a flattened joined table, as structured
relational input, or as a graph.
Yet no benchmark enforces identical task semantics across all three
representations simultaneously, making it impossible to determine whether
observed performance differences reflect inductive bias or simply
differences in available information.

We introduce \cb{}, a tri-view evaluation framework that holds the
prediction task and underlying data fixed while varying only the
representation.
Given a relational database and a temporally safe prediction task, \cb{}
derives three semantically equivalent views---a joined-table view
(\viewname{JT-Entity}, \viewname{JT-TemporalAgg}), a schema-native
relational view, and a graphified view---and evaluates models under a
unified protocol with leakage-safe, seed-time semantics.

We instantiate \cb{} on \todo{N} tasks drawn from
\todo{M} RelBench datasets and report results for tabular foundation
models (TabPFN, TabICL) on joined-table views.
Our experiments reveal three consistent findings: (i)~static entity
snapshots lose critical historical signal and perform near or below
chance; (ii)~temporal aggregate features recover most of this signal,
with longer lookback windows monotonically improving performance;
(iii)~representation sensitivity, as measured by our
\emph{Representation Invariance Score} (RIS), varies substantially
across tasks, indicating that no single representation universally
dominates.
We release the \cb{} protocol, data-construction pipeline, and
evaluation harness to support reproducible comparison across the
relational learning landscape.
\end{abstract}

% -----------------------------------------------------------------------
\section{Introduction}
\label{sec:intro}
% -----------------------------------------------------------------------

Relational databases are the dominant storage format for structured
enterprise and scientific data.
Training machine learning models on such data requires choosing a
\emph{representation}: practitioners either (i) denormalize the database
into a single flat table via feature engineering, (ii) feed normalized
tables directly into schema-native relational models, or (iii) convert
the database into a graph and apply graph neural networks.

Each choice embeds a different inductive bias.
Joined-table approaches are flexible but can lose inter-row dependencies
and are sensitive to aggregation design~\citep{relbench}.
Schema-native models such as the Relational Transformer operate directly
on multi-table inputs and may be more robust to feature engineering
choices~\citep{rt}.
Graph-based models such as RelGT can capture structural connectivity but
depend critically on how the graph is constructed~\citep{relgt}.

\paragraph{The problem.}
Despite this diversity of approaches, rigorous head-to-head comparison
is difficult because existing evaluations change both the representation
\emph{and} the information available to the model simultaneously.
A model given a carefully engineered feature matrix may outperform a
model given raw normalized tables for reasons that have nothing to do
with inductive bias---the information content of the two inputs simply
differs.
We formalize this as the \textbf{representation invariance problem}:
given a fixed relational database and prediction task, how sensitive is
model performance to the chosen representation when information content
is held constant?

\paragraph{Our contribution.}
We introduce \cb{}, a tri-view evaluation protocol designed to
make this question empirically testable.
\cb{} derives multiple semantically equivalent \emph{views} from the
same underlying data and evaluates models under a unified protocol with
strict temporal safety guarantees.
Our specific contributions are:

\begin{enumerate}
  \item \textbf{Tri-view protocol.} We define four views---\viewname{JT-Entity},
    \viewname{JT-TemporalAgg}, a schema-native relational view, and a
    graphified view---that preserve identical task semantics and temporal
    safety across representations (\S\ref{sec:protocol}).

  \item \textbf{Leakage-safe data construction pipeline.}
    We implement backward as-of joins and lookback-window aggregation
    that ensure no post-seed information is used in any view
    (\S\ref{sec:pipeline}).

  \item \textbf{Representation Invariance Score (RIS).}
    We propose a formal metric for cross-view consistency that
    decomposes variance in performance into representation-driven and
    task-driven components (\S\ref{sec:ris}).

  \item \textbf{Empirical evaluation.}
    We report Protocol~A (cross-view comparison) and Protocol~B
    (brittleness analysis) results on \todo{N tasks from M datasets},
    with \todo{full RT and RelGT baselines to be added}
    (\S\ref{sec:experiments}).
\end{enumerate}

% -----------------------------------------------------------------------
\section{Related Work}
\label{sec:related}
% -----------------------------------------------------------------------

\paragraph{Relational machine learning benchmarks.}
RelBench~\citep{relbench} is the closest prior work: it provides
temporally safe prediction tasks on multi-table relational databases
and supports tabular, relational, and graph-based baselines.
\cb{} builds directly on RelBench tasks and temporal splits, but adds
the tri-view protocol that holds task semantics fixed across
representations.
Without this control, differences in RelBench results across model
families confound representation choice with feature engineering
decisions.

\paragraph{Automated feature engineering.}
Systems such as FeatureTools and OpenFE automate joined-table
construction from relational databases, but do not study how
representation choice affects downstream performance~\citep{featuretools}.
\cb{} treats feature engineering choices as experimental variables
rather than implementation details.

\paragraph{Tabular foundation models.}
TabPFN~\citep{tabpfn} and TabICL~\citep{tabicl} are in-context learners
for small tabular datasets.
We use them as controlled baselines on the joined-table views, where
their closed-form inference makes experimental turnaround fast.

\paragraph{Schema-native relational models.}
The Relational Transformer~\citep{rt} encodes multi-table relational
structure natively, without requiring a flat-table preprocessing step.
We target it as the baseline for the schema-native view in \cb{}.

\paragraph{Relational graph models.}
RelGT~\citep{relgt} converts relational databases into heterogeneous
graphs and applies a graph transformer.
Its sensitivity to graph construction parameters (neighborhood size $K$)
makes it a natural candidate for Protocol~B brittleness analysis.

\paragraph{Multi-view learning.}
The multi-view learning literature studies fusion of heterogeneous
input modalities~\citep{multiview_survey}.
\cb{} differs in that the views are not independent sources of
information---they are derived from the same data---and the goal is
evaluation rather than fusion.

% -----------------------------------------------------------------------
\section{The \cb{} Protocol}
\label{sec:protocol}
% -----------------------------------------------------------------------

\subsection{Problem formulation}

Let $\mathcal{D}$ be a relational database and $\mathcal{T}$ a
prediction task defined by an entity table, a target column, and a
set of prediction instances
\[
  \{(e_i,\, t_i,\, y_i)\}_{i=1}^{N},
\]
where $e_i$ is the entity identifier, $t_i \in \mathbb{R}$ is the
seed time, and $y_i$ is the target label or value.

A \emph{representation} is a function $\phi$ that maps
$(\mathcal{D}, e_i, t_i)$ to a feature vector or structured input
consumed by a model $f_\phi$.
The \textbf{representation invariance problem} asks: for a fixed
$(\mathcal{D}, \mathcal{T})$, how much does
$\mathbb{E}[\ell(f_\phi, y)]$ vary across valid choices of $\phi$?

\subsection{Temporal safety}

All views enforce \emph{seed-time semantics}: only information with
timestamp $\tau \le t_i$ is available to the model.
For a record table $R$ with timestamps $\tau$, the leakage-safe
backward as-of join selects

\[
  r^\star = \arg\max_{r \in R(e_i)} \{\tau_r \mid \tau_r \le t_i\},
\]

and the lookback-window restriction is
\[
  R_W(e_i, t_i) = \{r \in R(e_i) \mid t_i - W \le \tau_r \le t_i\}.
\]

Temporal train/validation/test splits follow RelBench conventions
so that no future label information contaminates model selection.

\subsection{The four views}
\label{sec:views}

\paragraph{View 1: \viewname{JT-Entity}.}
Entity attributes are denormalized into the task table via leakage-safe
as-of joins (or static joins for tables without a time column).
This is the standard joined-table baseline used in tabular ML.
It contains no summarized event history.

\paragraph{View 2: \viewname{JT-TemporalAgg}.}
Starting from \viewname{JT-Entity}, we add leakage-safe aggregate
features computed over $R_W(e_i, t_i)$: row count, recency (days since
last event), mean, sum, min, max, and number of unique values of each
numerical column.
The lookback window $W$ is a design parameter studied in Protocol~B.

\paragraph{View 3: Relational (schema-native).}
The normalized tables and task metadata are preserved without
flattening:
\begin{lstlisting}
{
  "db": DatabaseSpec(tables={"entity_table": ...,
                             "event_table":  ...}),
  "task": TaskSpec(...)
}
\end{lstlisting}
This view is consumed by schema-native models such as the Relational
Transformer~\citep{rt}.

\paragraph{View 4: Structural-count proxy (current graph track).}
The relational payload is augmented with graph construction metadata:
\begin{lstlisting}
{
  "db": DatabaseSpec(...),
  "task": TaskSpec(...),
  "graph_config": {"neighborhood_size_k": K}
}
\end{lstlisting}
In the current benchmark runs, this track is materialized as a
count-feature proxy: for each target entity at seed time, it computes
incident row-count features from linked tables (bounded by $K$).
This makes the current graph track a structural tabular proxy rather than
a full per-instance \texttt{HeteroData} graph pipeline.
A full graphified implementation for GNN/graph-transformer models remains
part of the next implementation stage.

Table~\ref{tab:view_summary} summarizes the four views and their
target model families.

\begin{table}[t]
\centering
\caption{Summary of \cb{} views. All views share identical task
semantics and temporal safety guarantees.}
\label{tab:view_summary}
\begin{tabular}{lllp{5cm}}
\toprule
View & Format & History & Target models \\
\midrule
\viewname{JT-Entity}       & Flat table   & None (static)         & TabPFN, TabICL, XGBoost \\
\viewname{JT-TemporalAgg}  & Flat table   & Aggregated ($W$ days) & TabPFN, TabICL, XGBoost \\
Relational                 & Multi-table  & Full (normalized)     & Relational Transformer (subprocess) \\
Structural-count proxy     & Flat table (count features)  & Structural counts ($\le K$ linked rows/table) & TabPFN, TabICL, XGBoost (current); RGCN/GraphSAGE/RelGT/ULTRA in planned full graph view \\
\bottomrule
\end{tabular}
\end{table}

\subsection{Evaluation protocols}

\paragraph{Protocol A: Cross-view comparison.}
We compare all four views on identical prediction instances, using the
best-performing model for each view family.
Metrics include predictive performance (AUROC for classification,
MAE for regression), compute cost (GPU-hours, peak VRAM), and RIS
(Section~\ref{sec:ris}).

\paragraph{Protocol B: Brittleness analysis.}
We ablate each view's construction parameters:
\begin{itemize}
  \item lookback window $W \in \{7, 30, 90, \mathrm{None}\}$ for
        \viewname{JT-TemporalAgg},
  \item join depth $d \in \{1, 2\}$ for \viewname{JT-Entity},
  \item graph neighborhood budget $K \in \{100, 300, 500\}$ for the
        structural-count proxy track.
\end{itemize}
High sensitivity to these parameters indicates brittleness in that
representation family.

% -----------------------------------------------------------------------
\section{Representation Invariance Score}
\label{sec:ris}
% -----------------------------------------------------------------------

Let $\mathcal{V} = \{v_1, \ldots, v_m\}$ be the set of views and let
$\mu_v$ denote the evaluation metric (e.g., AUROC) achieved by the
best model in view $v$.
To make scores comparable across tasks with different baseline
difficulties, we normalize each view's score to a $[0,1]$ utility:
\[
  u_v = \frac{\mu_v - \mu_{\text{chance}}}{\mu_{\text{oracle}} - \mu_{\text{chance}}},
\]
where $\mu_{\text{chance}}$ is chance-level performance (0.5 for AUROC)
and $\mu_{\text{oracle}}$ is the best score across all views and models
on that task (an empirical oracle ceiling).

The \textbf{Representation Invariance Score} for task $\mathcal{T}$ is:
\[
  \RIS(\mathcal{T}) = 1 - \sigma(\{u_v\}_{v \in \mathcal{V}}),
\]
where $\sigma(\cdot)$ denotes the standard deviation across views.
$\RIS \in [0, 1]$, where $\RIS = 1$ indicates that all representations
achieve equal normalized utility (perfect invariance) and lower values
indicate greater representation sensitivity.

\paragraph{Interpretation.}
A task with $\RIS < 0.7$ is considered \emph{representation-sensitive}:
the choice of representation has a material impact on achievable
performance.
A task with $\RIS \ge 0.9$ is \emph{representation-robust}: any of the
three representation families suffices.
\todo{Calibrate these thresholds empirically across all benchmark
tasks.}

\paragraph{Decomposition.}
Across a suite of tasks $\{\mathcal{T}_j\}$, we also report the
\emph{macro-RIS}:
\[
  \overline{\RIS} = \frac{1}{|\mathcal{T}|}\sum_j \RIS(\mathcal{T}_j),
\]
and its standard deviation, to summarize benchmark-level consistency.

% -----------------------------------------------------------------------
\section{Data Construction Pipeline}
\label{sec:pipeline}
% -----------------------------------------------------------------------

We implement the \cb{} data construction pipeline on top of RelBench.
The pipeline takes as input a RelBench dataset and task specification
and outputs all four views in a unified format.

\subsection{Input: RelBench relational format}

The raw data is stored in normalized relational form.
For example, for the \texttt{rel-f1 / driver-top3} task, each
prediction instance is a row \texttt{(date, driverId, qualifying)}
in the task table.
To recover context, one must join against \texttt{drivers},
\texttt{qualifying}, \texttt{results}, and \texttt{standings}
tables (Tables~\ref{tab:relbench-task-example}--\ref{tab:relbench-history-example}).

\begin{table}[h]
\centering
\small
\caption{Example task row in RelBench for \texttt{driver-top3}.}
\label{tab:relbench-task-example}
\begin{tabular}{lll}
\toprule
date & driverId & qualifying \\
\midrule
2004-08-04 & 12 & 0 \\
\bottomrule
\end{tabular}
\end{table}

\begin{table}[h]
\centering
\small
\caption{Corresponding row from the raw \texttt{drivers} table.}
\label{tab:relbench-drivers-example}
\begin{tabular}{lllllll}
\toprule
driverId & driverRef & code & forename & surname & dob & nationality \\
\midrule
12 & massa & MAS & Felipe & Massa & 1981-04-25 & Brazilian \\
\bottomrule
\end{tabular}
\end{table}

\begin{table}[h]
\centering
\small
\caption{Historical rows from \texttt{qualifying} and \texttt{results}
linked to the same driver, available before the seed time.}
\label{tab:relbench-history-example}
\begin{tabular}{lllllll}
\toprule
table & date & raceId & driverId & position & extra \\
\midrule
qualifying & 2004-07-24 & 724 & 12 & 16 & qualifyId=2091 \\
qualifying & 2004-07-10 & 723 & 12 & 11 & qualifyId=2068 \\
results    & 2004-07-25 & 724 & 12 & 13 & grid=16, points=0 \\
results    & 2004-07-11 & 723 & 12 & 9  & grid=10, points=0 \\
\bottomrule
\end{tabular}
\end{table}

\subsection{View construction}

\paragraph{\viewname{JT-Entity}.}
We apply a leakage-safe as-of join from the task table to each
related entity table.
For tables without a time column, a static join is used.
The result is one flat row per prediction instance containing
entity-level attributes only (no event history).

\paragraph{\viewname{JT-TemporalAgg}.}
We extend \viewname{JT-Entity} by computing, for each related event
table $R$ and each numerical column $c$, the following statistics
over $R_W(e_i, t_i)$:
\begin{align*}
  &\texttt{count},\ \texttt{recency\_days},\ \texttt{mean}(c),\
   \texttt{sum}(c),\ \texttt{min}(c),\ \texttt{max}(c),\ \texttt{nunique}(c).
\end{align*}
Missing values arise when no events fall within the window.

\paragraph{Relational and structural-count proxy views.}
Both views wrap the normalized database and task specification in a
structured payload, with the structural-count proxy adding
\texttt{neighborhood\_size\_k} as a construction parameter.

\subsection{Actual materialized examples for each view}
\label{sec:actual_view_examples}

To make the view differences concrete, we show actual rows from the
\texttt{rel-f1/driver-top3} task for the same prediction instance:
\texttt{(date = 2004-08-04, driverId = 12, qualifying = 0)}.

\paragraph{\viewname{JT-Entity} (actual row).}
\begin{table}[h]
\centering
\small
\caption{Actual \viewname{JT-Entity} row (selected columns).}
\label{tab:jt_entity_actual_example}
\begin{tabular}{llllllll}
\toprule
date & driverId & qualifying & driverRef & code & forename & surname & nationality \\
\midrule
2004-08-04 & 12 & 0 & massa & MAS & Felipe & Massa & Brazilian \\
\bottomrule
\end{tabular}
\end{table}

\paragraph{\viewname{JT-TemporalAgg} (actual row).}
\begin{table}[h]
\centering
\small
\caption{Actual \viewname{JT-TemporalAgg} row (selected aggregate columns).}
\label{tab:jt_temporalagg_actual_example}
\begin{tabular}{llllllllll}
\toprule
date & driverId & qualifying & qual\_cnt & qual\_recency & qual\_pos\_mean & res\_cnt & res\_pos\_mean & std\_cnt & std\_pos\_mean \\
\midrule
2004-08-04 & 12 & 0 & 2 & 11.0 & 13.5 & 2 & 11.0 & 2 & 12.5 \\
\bottomrule
\end{tabular}
\end{table}

In Table~\ref{tab:jt_temporalagg_actual_example}, the shorthand columns map to:
\texttt{qual\_cnt = qualifying\_\_row\_count},
\texttt{qual\_recency = qualifying\_\_recency\_days},
\texttt{qual\_pos\_mean = qualifying\_\_position\_\_mean},
\texttt{res\_cnt = results\_\_row\_count},
\texttt{res\_pos\_mean = results\_\_position\_\_mean},
\texttt{std\_cnt = standings\_\_row\_count},
\texttt{std\_pos\_mean = standings\_\_position\_\_mean}.

\paragraph{Relational view (actual payload shape).}
\begin{lstlisting}
{
  "db": DatabaseSpec(tables={
      "drivers": ...,
      "qualifying": ...,
      "results": ...,
      "standings": ...,
      ...
  }),
  "task": TaskSpec(
      name="rel-f1__driver-top3",
      task_type="classification",
      target_col="qualifying",
      entity_table="drivers",
      entity_key="driverId",
      seed_time_col="date"
  )
}
\end{lstlisting}

\paragraph{Structural-count proxy view (actual payload and derived row).}
For the current benchmark runs in this repository, the graph branch is a
structural count-based proxy. The payload and resulting row for the same
instance are:
\begin{lstlisting}
{
  "db": DatabaseSpec(...),
  "task": TaskSpec(...),
  "graph_config": {"neighborhood_size_k": 300}
}
\end{lstlisting}

\begin{table}[h]
\centering
\small
\caption{Actual structural-count proxy row for the same instance.}
\label{tab:graphified_actual_example}
\begin{tabular}{llllll}
\toprule
date & driverId & qualifying & qual\_count\_leq\_seed & result\_count\_leq\_seed & standing\_count\_leq\_seed \\
\midrule
2004-08-04 & 12 & 0 & 14 & 28 & 27 \\
\bottomrule
\end{tabular}
\end{table}

The count columns in Table~\ref{tab:graphified_actual_example} correspond to:
\texttt{graphdeg\_\_qualifying\_\_count\_leq\_seed = 14},
\texttt{graphdeg\_\_results\_\_count\_leq\_seed = 28},
\texttt{graphdeg\_\_standings\_\_count\_leq\_seed = 27}.
These values are exactly the number of raw rows linked to
\texttt{driverId=12} in each table with timestamp \(\le\) seed time.

\subsection{Implemented validation}

We validated the full pipeline on a synthetic benchmark with a
\texttt{users}/\texttt{transactions} schema.
The following components have been verified:
\begin{itemize}
  \item leakage-safe backward as-of joins produce no post-seed
        information,
  \item temporal splits are consistent across all four views,
  \item \viewname{JT-Entity} and \viewname{JT-TemporalAgg} construction
        is correct on the synthetic task,
  \item relational and structural-count proxy payloads are well-formed and
        accepted by downstream model interfaces.
\end{itemize}

% -----------------------------------------------------------------------
\section{Experiments}
\label{sec:experiments}
% -----------------------------------------------------------------------

\subsection{Datasets and tasks}

\paragraph{Current results: \texttt{rel-f1 / driver-top3}.}
This task asks whether a given Formula~1 driver will qualify in the
top three in a race within the next month.
It is a binary classification task evaluated with AUROC.
We use the official RelBench temporal split.

\paragraph{\todo{Additional tasks (to be added before submission).}}
We plan to evaluate \cb{} on the following additional RelBench tasks
to ensure generality of findings:
\begin{itemize}
  \item \todo{\texttt{rel-amazon / user-churn} (e-commerce, regression)},
  \item \todo{\texttt{rel-stack / post-votes} (community Q\&A, regression)},
  \item \todo{\texttt{rel-trial / site-success} (clinical trial, classification)},
  \item \todo{at least one additional classification task from a
        different domain to cover at least 4 tasks across 3 datasets}.
\end{itemize}
This coverage ensures that conclusions about representation
sensitivity are not specific to the Formula~1 domain.

\subsection{Models and baselines}

For joined-table views, we use:
\begin{itemize}
  \item \textbf{TabPFN}~\citep{tabpfn}: a prior-fitted Bayesian
        in-context learner for small tabular datasets.
  \item \textbf{TabICL}~\citep{tabicl}: a larger in-context learner
        with improved scaling.
  \item \textbf{XGBoost}: a gradient-boosted tree baseline
        (\texttt{XGBClassifier}/\texttt{XGBRegressor}, \texttt{tree\_method=hist},
        500 estimators) for comparison with foundation models.
\end{itemize}

For the relational view:
\begin{itemize}
  \item \textbf{Relational Transformer}~\citep{rt}: the schema-native
        baseline, integrated via a subprocess adapter
        (\texttt{OfficialRelationalTransformerAdapter}) that invokes
        the official RT implementation and scrapes the reported test metric.
        Requires Linux and CUDA; currently returns a \texttt{precomputed\_test\_metric}.
\end{itemize}

For the graph track:
\begin{itemize}
  \item \textbf{RGCN}: a relational graph convolutional network
        (\texttt{GCNConv} with per-type projections) trained in-process
        via PyTorch Geometric.
  \item \textbf{GraphSAGE}: a heterogeneous GraphSAGE
        (\texttt{SAGEConv} per edge type) trained in-process.
  \item \textbf{RelGT}~\citep{relgt}: a graph transformer for
        relational databases, integrated via a subprocess adapter
        (\texttt{RelGTSubprocessAdapter}); requires Linux and CUDA.
  \item \textbf{ULTRA}: a universal link-prediction foundation model,
        integrated via a subprocess adapter (\texttt{ULTRASubprocessAdapter}).
        Note: ULTRA is designed for KG link prediction; applying it to
        entity-property prediction requires reformulation as link scoring.
\end{itemize}

\paragraph{Implementation status.}
RGCN and GraphSAGE are fully in-process and produce row-level predictions.
Both share a common training loop with Adam, \texttt{ReduceLROnPlateau},
early stopping, class-balance weighting (\texttt{pos\_weight}), and
gradient clipping.
RelGT, ULTRA, and RT are integrated via subprocess adapters and return
pre-computed test metrics.
Results in Tables~\ref{tab:protocol-a-driver-top3}
and~\ref{tab:protocol-b-driver-top3} correspond to the
TabPFN run (job \texttt{21782164}), TabICL run (job \texttt{21782181}),
and official RT run (job \texttt{21786870}).
The official RT run has been completed, achieving AUROC~0.788788 on
\texttt{rel-f1/driver-top3}, which is the highest Protocol~A score
across all models and views for this task.
Full graph-model runs (RGCN/GraphSAGE/RelGT/ULTRA) are
\todo{pending execution}.
GPU peak VRAM is tracked automatically via
\texttt{torch.cuda.max\_memory\_allocated()} for all GPU-accelerated
models and reported in the compute-cost analysis.

\subsection{Protocol A results}
\label{sec:protocol_a_results}

\begin{table}[t]
\centering
\caption{Protocol~A results on \texttt{rel-f1/driver-top3}.
         AUROC; chance level = 0.5. TabPFN/TabICL joined-table views
         from Snellius jobs \texttt{21782164} and \texttt{21782181};
         RT-official from job \texttt{21786870}.}
\label{tab:protocol-a-driver-top3}
\begin{tabular}{llll}
\toprule
Model & View & AUROC & Notes \\
\midrule
TabPFN & \viewname{JT-Entity}                      & 0.469142 & joined-table \\
TabPFN & \viewname{JT-TemporalAgg} ($W=30$, $d=1$) & 0.682405 & joined-table \\
TabPFN & Structural-count proxy ($K=300$)          & 0.669726 & graph track proxy \\
TabICL & \viewname{JT-Entity}                      & 0.758720 & joined-table \\
TabICL & \viewname{JT-TemporalAgg} ($W=30$, $d=1$) & 0.509530 & joined-table \\
TabICL & Structural-count proxy ($K=300$)          & 0.611446 & graph track proxy \\
\midrule
RT-official & Relational                           & \textbf{0.788788} & job 21786870 \\
\bottomrule
\end{tabular}
\end{table}

Table~\ref{tab:protocol-a-driver-top3} shows model-dependent behavior.
For TabPFN, \viewname{JT-Entity} is below chance (0.469142), while
\viewname{JT-TemporalAgg} (0.682405) and structural-count proxy
(0.669726) are substantially stronger.
For TabICL, the pattern differs: \viewname{JT-Entity} is strongest
(0.758720), \viewname{JT-TemporalAgg} at 30 days drops near chance
(0.509530), and structural-count proxy is intermediate (0.611446).
The official Relational Transformer achieves 0.788788, the highest
Protocol~A score across all views and models on this task, exceeding
the best joined-table result (TabICL \viewname{JT-Entity}: 0.758720)
by 3.0 AUROC points.

\paragraph{Protocol A RIS.}
The harness reports per-model RIS (computed over the views each model was evaluated on):
\begin{itemize}
  \item TabPFN (3 views: JT-Entity, JT-TemporalAgg $W=30$, structural proxy):
        \(\RIS = 0.471116\),
        \(\texttt{utility\_std} = 0.528884\),
        \(\texttt{utility\_gap} = 1.169173\).
  \item TabICL (3 views): \(\RIS = 0.597368\),
        \(\texttt{utility\_std} = 0.402632\),
        \(\texttt{utility\_gap} = 1.000000\).
  \item RT-official (1 view: relational): \(\RIS = 1.000000\),
        \(\texttt{utility\_std} = 0.000000\),
        \(\texttt{utility\_gap} = 0.000000\)
        (trivially 1.0 when only one view is evaluated).
\end{itemize}
Both TabPFN and TabICL values indicate non-trivial representation sensitivity.
A joint cross-model Protocol~A RIS incorporating RT-official in the oracle
ceiling is deferred to the full multi-task analysis.

\paragraph{Contrasting entity behavior.}
The 0.469142 AUROC of TabPFN on \viewname{JT-Entity} merits a note.
At a \emph{dataset level}, chance performance is 0.5; individual
model runs can fall slightly below due to finite-sample variance,
class imbalance, or overfitting on small splits.
In contrast, TabICL reaches 0.758720 on the same view, indicating that
representation effects are strongly model-dependent even under identical
task semantics.
\todo{Report class imbalance ratio and confirm with multiple random
seeds.}

\subsection{Protocol B ablation results}
\label{sec:protocol_b_results}

\begin{table}[t]
\centering
\caption{Protocol~B ablations on \texttt{rel-f1/driver-top3} with TabPFN
         (job \texttt{21782164}).}
\label{tab:protocol-b-driver-top3}
\begin{tabular}{llll}
\toprule
Parameter & Setting & View & AUROC \\
\midrule
\multirow{4}{*}{Lookback $W$}
  & 7 days          & \viewname{JT-TemporalAgg} ($d=1$) & 0.611400 \\
  & 30 days         & \viewname{JT-TemporalAgg} ($d=1$) & 0.682405 \\
  & 90 days         & \viewname{JT-TemporalAgg} ($d=1$) & 0.732056 \\
  & $\infty$ (all history) & \viewname{JT-TemporalAgg} ($d=1$) & \textbf{0.802284} \\
\midrule
\multirow{2}{*}{Join depth $d$}
  & $d=1$ & \viewname{JT-Entity} ($W=\mathrm{None}$) & 0.469142 \\
  & $d=2$ & \viewname{JT-Entity} ($W=\mathrm{None}$) & 0.469142 \\
\midrule
\multirow{3}{*}{Neighborhood $K$}
  & 100 & Structural-count proxy & 0.576381 \\
  & 300 & Structural-count proxy & 0.669726 \\
  & 500 & Structural-count proxy & 0.669726 \\
\bottomrule
\end{tabular}
\end{table}

\begin{table}[t]
\centering
\caption{Protocol~B ablations on \texttt{rel-f1/driver-top3} with TabICL
         (job \texttt{21782181}).}
\label{tab:protocol-b-driver-top3-tabicl}
\begin{tabular}{llll}
\toprule
Parameter & Setting & View & AUROC \\
\midrule
\multirow{4}{*}{Lookback $W$}
  & 7 days          & \viewname{JT-TemporalAgg} ($d=1$) & 0.588485 \\
  & 30 days         & \viewname{JT-TemporalAgg} ($d=1$) & 0.510125 \\
  & 90 days         & \viewname{JT-TemporalAgg} ($d=1$) & 0.580738 \\
  & $\infty$ (all history) & \viewname{JT-TemporalAgg} ($d=1$) & \textbf{0.759909} \\
\midrule
\multirow{2}{*}{Join depth $d$}
  & $d=1$ & \viewname{JT-Entity} ($W=\mathrm{None}$) & 0.758720 \\
  & $d=2$ & \viewname{JT-Entity} ($W=\mathrm{None}$) & 0.758720 \\
\midrule
\multirow{3}{*}{Neighborhood $K$}
  & 100 & Structural-count proxy & 0.566021 \\
  & 300 & Structural-count proxy & 0.611446 \\
  & 500 & Structural-count proxy & 0.611446 \\
\bottomrule
\end{tabular}
\end{table}

Protocol~B results in Tables~\ref{tab:protocol-b-driver-top3}
and~\ref{tab:protocol-b-driver-top3-tabicl} clarify the source of
performance variation.

\paragraph{Window sensitivity (\viewname{JT-TemporalAgg}) is model-specific.}
For TabPFN, performance rises from 0.611400 ($W=7$) to 0.802284
($W=\infty$), indicating a strong long-history benefit.
For TabICL, $W=30$ is weak (0.510125), while $W=\infty$ recovers to
0.759909, closer to the entity view. This suggests the effect of
aggregation window depends on model family, not only on task signal.

\paragraph{Join-depth sensitivity (\viewname{JT-Entity}).}
We ablate join depth $d \in \{1, 2\}$.
At $d=1$, only direct entity attributes are included (e.g., the
\texttt{drivers} row).
At $d=2$, attributes of entities reachable via one FK hop are also
denormalized (e.g., constructor attributes).
For this run, $d=2$ matches $d=1$ exactly (both 0.469142), suggesting
that deeper static joins do not help for TabPFN.
For TabICL, $d=2$ also matches $d=1$ exactly (both 0.758720).

\paragraph{Neighborhood sensitivity (structural-count proxy).}
Increasing $K$ from 100 to 300 improves AUROC for both models
(TabPFN: 0.576381 \(\to\) 0.669726;
TabICL: 0.566021 \(\to\) 0.611446),
but further increasing to $K=500$ yields no improvement in either run.
This saturation is consistent with the small number of incident tables
in the \texttt{rel-f1} schema: once the neighborhood budget covers
all linked qualifying and results records, additional budget adds
no new information.
\todo{Verify with a schema that has deeper relational structure
(e.g., \texttt{rel-amazon}) to test whether saturation persists.}

\paragraph{RIS computation (latest run).}
The harness reports the following RIS values for
\texttt{rel-f1/driver-top3}:
\begin{itemize}
  \item TabPFN, Protocol~A: \(\RIS = 0.471116\),
        \(\texttt{utility\_std} = 0.528884\),
        \(\texttt{utility\_gap} = 1.169173\).
  \item TabPFN, Protocol~B: \(\RIS = 0.687582\),
        \(\texttt{utility\_std} = 0.312418\),
        \(\texttt{utility\_gap} = 1.102083\).
  \item TabICL, Protocol~A: \(\RIS = 0.597368\),
        \(\texttt{utility\_std} = 0.402632\),
        \(\texttt{utility\_gap} = 1.000000\).
  \item TabICL, Protocol~B: \(\RIS = 0.676962\),
        \(\texttt{utility\_std} = 0.323038\),
        \(\texttt{utility\_gap} = 0.961045\).
\end{itemize}
Across both models, Protocol~A and Protocol~B indicate non-trivial
representation sensitivity.
Macro-RIS across the full task suite will be reported via
\texttt{compute\_macro\_ris()}, which averages per-task RIS values and
reports the standard deviation across tasks.

\subsection{\todo{Planned experiments}}

The following experiments are required before submission to establish
the generality of \cb{} findings:

\begin{enumerate}
  \item \textbf{Additional tasks.} Run Protocol~A and B on the four
        additional RelBench tasks listed in Section~\ref{sec:experiments}.
        Report per-task AUROC/MAE and RIS, and macro-RIS across tasks
        via \texttt{compute\_macro\_ris()}.
  \item \textbf{Full relational baseline.} Execute the official
        Relational Transformer via \texttt{OfficialRelationalTransformerAdapter}
        (subprocess) on a Linux+CUDA node for all tasks.
        [\emph{Completed for \texttt{rel-f1/driver-top3}:
        job \texttt{21786870}, AUROC = 0.788788.
        Remaining tasks pending.}]
  \item \textbf{In-process GNN baselines.} Execute RGCN and GraphSAGE
        (fully implemented in \texttt{adapter\_graph\_nn.py}) on the
        full graphified view for all tasks.
        [\emph{Models implemented; execution pending.}]
  \item \textbf{RelGT and ULTRA baselines.} Execute RelGT via
        \texttt{RelGTSubprocessAdapter} and ULTRA via
        \texttt{ULTRASubprocessAdapter} on a Linux+CUDA node.
        [\emph{Adapters implemented; execution pending. Note: ULTRA
        requires reformulation from KG link-prediction to entity
        classification.}]
  \item \textbf{XGBoost baseline.} Execute XGBoost
        (\texttt{build\_xgboost()} via \texttt{--models xgboost})
        on \viewname{JT-Entity} and \viewname{JT-TemporalAgg}.
        [\emph{Implemented; execution pending.}]
  \item \textbf{Protocol~B join-depth and 90-day ablations.}
        Run $d=2$ (\viewname{JT-Entity}) and $W=90$ (\viewname{JT-TemporalAgg})
        ablations now wired into the default Protocol~B configuration.
        [\emph{Configuration implemented; execution pending.}]
  \item \textbf{Compute cost analysis.} Report GPU-hours and peak VRAM
        (tracked via \texttt{peak\_gpu\_vram\_mb} in \texttt{compute\_tracker})
        for each model/view combination.
  \item \textbf{Multi-seed evaluation.} Report mean $\pm$ std across
        at least 3 random seeds to confirm stability of AUROC
        estimates.
\end{enumerate}

% -----------------------------------------------------------------------
\section{Discussion}
\label{sec:discussion}
% -----------------------------------------------------------------------

\paragraph{Representation sensitivity is task-dependent.}
The \texttt{driver-top3} experiments show that representation choice
has a large effect on achievable performance for this task.
However, the preferred view depends on the model: TabPFN benefits from
temporal aggregation, while TabICL is strongest on the static entity
view in the current runs.
Whether this pattern generalizes to other relational tasks---or whether
some tasks are truly representation-invariant---is precisely the
question \cb{} is designed to answer at scale.

\paragraph{Temporal depth and model interaction.}
For TabPFN, extending lookback from 30 days to full history yields a
large gain (0.682405 $\to$ 0.802284). For TabICL, full-history
aggregation also helps relative to 30 days
(0.510125 $\to$ 0.759909), but static entity remains competitive.
This indicates that temporal feature design and model choice interact
strongly, rather than one universally dominating the other.

\paragraph{Structural neighborhood saturation.}
The structural-count proxy saturates at $K = 300$ for \texttt{rel-f1}.
This is expected given the shallow schema structure of the dataset.
Richer schemas with deeper relational hierarchies may require larger
budgets and may show greater variance in performance with $K$, making
Protocol~B a useful diagnostic tool.

\paragraph{Limitations.}
The current results cover a single dataset and task.
The relational track now has an official RT result (AUROC~0.788788,
job \texttt{21786870}), but the graph track still uses a
structural-count proxy rather than a full GNN.
All conclusions should therefore be treated as preliminary until the
full experimental suite across all tasks is complete.
Additionally, RIS as currently defined uses an empirical oracle
ceiling, which means it cannot be compared across benchmarks without
a common normalization scheme.
\todo{Consider fixed, task-type-specific oracle ceilings based on
upper-bound estimates from the literature.}

% -----------------------------------------------------------------------
\section{Conclusion}
\label{sec:conclusion}
% -----------------------------------------------------------------------

We have presented \cb{}, a tri-view benchmark for studying
representation invariance in relational machine learning.
\cb{} enforces identical task semantics and temporal safety guarantees
across four derived views---joined-table entity, joined-table temporal
aggregate, schema-native relational, and a graph-oriented
structural-count proxy---and provides a
unified evaluation protocol for comparing tabular, relational, and
graph-based models.

Our initial results on \texttt{rel-f1 / driver-top3} demonstrate that
representation choice has a substantial impact on predictive performance
for this task.
The official Relational Transformer achieves AUROC~0.788788, the highest
score across all Protocol~A views and models, exceeding the best
joined-table result (TabICL \viewname{JT-Entity}: 0.758720) by 3.0 points.
Among joined-table models, TabPFN is strongest with temporal aggregation,
while TabICL is strongest on the static entity view.
The Representation Invariance Score formalizes this sensitivity in a
task-agnostic way.

We are extending \cb{} to \todo{N} tasks across \todo{M} RelBench
datasets, and integrating full Relational Transformer and RelGT
baselines to enable true tri-view comparison.
We will release the full pipeline and evaluation harness to support
reproducible research on relational representation learning.

% -----------------------------------------------------------------------
\section*{Acknowledgements}
% -----------------------------------------------------------------------
\todo{Add acknowledgements, and compute credits.}

% -----------------------------------------------------------------------
\bibliographystyle{plainnat}
\begin{thebibliography}{9}

\bibitem[Fey et al.(2024)]{relbench}
Matthias Fey, Weihua Hu, Kexin Huang, Jan Eric Lenssen, Rishabh Ranjan,
Joshua Robinson, Rex Ying, Jiaxuan You, and Jure Leskovec.
\newblock RelBench: A benchmark for deep learning on relational databases.
\newblock In \emph{Advances in Neural Information Processing Systems
  (NeurIPS)}, 2024.

\bibitem[Ranjan et al.(2025)]{rt}
R.~Ranjan, V.~Hudovernik, M.~Znidar, C.~I.~Kanatsoulis,
R.~R.~Upendra, M.~Mohammadi, J.~Meyer, T.~Palczewski, C.~Guestrin,
and J.~Leskovec.
\newblock Relational Transformer: Toward zero-shot foundation models for
  relational data.
\newblock \emph{arXiv preprint arXiv:2510.06377}, 2025.
\newblock \textbf{[Note: verify this preprint exists and is citable before
  submission.]}

\bibitem[Dwivedi et al.(2025)]{relgt}
V.~P.~Dwivedi, S.~Jaladi, Y.~Shen, F.~L\'opez, C.~I.~Kanatsoulis,
R.~Puri, M.~Fey, and J.~Leskovec.
\newblock Relational Graph Transformer.
\newblock \emph{arXiv preprint arXiv:2505.10960}, 2025.

\bibitem[Hollmann et al.(2022)]{tabpfn}
Noah Hollmann, Samuel M\"uller, Katharina Eggensperger, and Frank Hutter.
\newblock TabPFN: A transformer that solves small tabular classification
  problems in a second.
\newblock In \emph{International Conference on Learning Representations
  (ICLR)}, 2023.

\bibitem[Ma et al.(2025)]{tabicl}
\todo{Add correct TabICL citation once publication details confirmed.}

\bibitem[Kanter and Veeramachaneni(2015)]{featuretools}
James Max Kanter and Kalyan Veeramachaneni.
\newblock Deep feature synthesis: Towards automating data science endeavors.
\newblock In \emph{IEEE International Conference on Data Science and Advanced
  Analytics}, 2015.

\bibitem[Xu et al.(2013)]{multiview_survey}
Chang Xu, Dacheng Tao, and Chao Xu.
\newblock A survey on multi-view learning.
\newblock \emph{arXiv preprint arXiv:1304.5634}, 2013.

\end{thebibliography}

\end{document}
